\section{Выводы}

В рамках лабораторной работы рассмотрена модель обмена Modbus с одним
ведущим и тремя ведомыми устройствами. На языке Promela описаны
адресная и широковещательная запись регистра, проверка получателя,
обработка неверной CRC, неподдерживаемой функции и ошибочного адреса
регистра. Для модели сформулированы и разобраны 14 LTL-свойств,
проверяющих сохранность памяти, корректность ответов и завершение обмена.

По результатам законченных поисков все свойства подтверждены,
недопустимых конечных состояний не обнаружено. При этом автоматический
запуск не смог закончить проверку свойства прогресса за установленное
время: его пришлось выполнить отдельно. Повторный поиск завершился
без нарушений за 214 секунд. Модель и формулы при этом не изменялись.

Существенную роль сыграли наблюдательные данные и согласование границ
транзакции. Снимки регистров, флаги операций, счетчик ответов и уведомления
обработчиков позволяют проверять именно результат рассматриваемого
запроса. Таким образом, при формализации системы необходимо учитывать
не только ее полезные действия, но и то, в какой момент сформулированное
свойство должно оцениваться.

Рассмотренная модель описывает только небольшой фрагмент поведения
Modbus. Даже при трех Slave, двух регистрах и ограниченном наборе
значений большинство проверок сохраняло более 14 миллионов состояний,
а проверка прогресса --- более 28 миллионов. Успешные поиски суммарно
заняли около 19 минут и потребовали до 2\,359 Мбайт памяти по отчету SPIN.
При расширении состава устройств, диапазонов данных и правил обмена
число сочетаний состояний и допустимых взаимодействий будет возрастать.
Поэтому верификация более полной системы потребует дальнейшего выбора
абстракций и, возможно, раздельной проверки ее частей. Полученный
результат показывает корректность выбранного фрагмента в пределах
явно заданных предположений.
